\vfill
\newPlot{wykres_Ct_Cn_eta_lambda.png}{Współczynniki}{wspolczynniki}
\akapit Na rysunku 1 widać, że współczynnik ciągu $C_T$ maleje ze wzrostem współczynnika posuwu $\lambda$, od ok. 0,145 dla $\lambda = 0$ do ok. 0,025 dla największego $\lambda$. Współczynnik momentu $C_N$ na początku jest prawie stały (ok. 0,01), a potem spada do zera. Sprawność śmigła dla małych $\lambda$ jest bliska zeru i rośnie dopiero na ostatnim punkcie.
\vfill\newpage
\newPlot{wykres_M_n}{Charakterystyka momentu}{moment}
\akapit Rysunek 2 pokazuje, że moment oporowy rośnie z prędkością obrotową mniej więcej kwadratowo, do ok. 8 Ncm. Wyniki dla wszystkich trzech warunków pomiaru prawie się pokrywają i są zgodne z modelem Quanser.

\vfill\newpage
\newPlot{wykres_I_n}{Charakterystyka prądu}{prunt}
\akapit Na rysunku 3 prąd rośnie z obrotami, od ok. 1 A do ok. 17 A. Nasze punkty leżą trochę powyżej krzywej katalogowej Cobra. Z tunelem prąd jest minimalnie większy przy dużych obrotach.

\vfill\newpage
\newPlot{wykres_T_n}{Charakterystyka siły ciągu}{ciag}
\akapit Siła ciągu (rysunek 4) rośnie z obrotami do ok. 7 N. Dla napięć 11,1 V i 9,5 V punkty leżą na jednej krzywej, czyli ciąg zależy od obrotów, a nie od napięcia. Z włączonym tunelem ciąg jest wyraźnie mniejszy (ok. 5,7 N przy największych obrotach). Wyniki są zbliżone do danych producenta (HQ 6x4.5) i modelu Quanser.

\vfill\newpage
\newPlot{wykres_sprawnosci}{Wykres sprawności}{sprawnosc}
\akapit Rysunek 5 pokazuje, że sprawność silnika $\eta_m$ rośnie z obrotami do ok. 0,7. Dla 9,5 V jest ona trochę większa przy małych obrotach. Sprawność całego zespołu napędowego $\eta_{zn}$ (z tunelem) ma maksimum ok. 0,35 przy ok. 175 obr/s, a potem lekko spada.
\vfill\newpage
